\section{A Day Each Month}

The \emph{Craft of Dying} complained that men will not ``dispose themselves to death betimes.'' The saints answered with a practice: to dispose oneself regularly, while in health, so that the disposition made at the end is one made many times before.

The simplest form of the practice is a choice made in the light of death. Ignatius of Loyola gives it as a rule for the election, the choice of a state or way of life: ``to consider, as if I were at the point of death, the form and measure which I would then want to have kept'' in the matter now to be chosen, and to decide accordingly.\endnote{Ignatius of Loyola, \emph{Spiritual Exercises}, Second Week, the second way to make a good and sound election, third rule, trans. Elder Mullan (New York: P. J. Kenedy \& Sons, 1914).} Francis de Sales ends his meditation on death with a resolution of the same kind: ``I will prepare for the hour of death and take every precaution for its peaceful arrival; I will thoroughly examine into the state of my conscience, and put in order whatever is wanting.''\endnote{Francis de Sales, \emph{Introduction to the Devout Life} I.13, Affections and Resolutions, Rivingtons.}

\subsection{The monthly retreat}

The fullest form is the monthly day of preparation for death, commended in the eighteenth century by Leonard of Port Maurice and printed with the English \emph{Preparation for Death} of Alphonsus. ``The great act of Christian prudence, in my opinion, consists in preparing to make well that last step, for which, if once made badly, there is no remedy.'' To be ready for it, Leonard says, two things are needed. The first is to pray constantly for the grace of final perseverance, ``which is the grace of graces,'' and to repeat often the words ``My Jesus, mercy, mercy,'' so that a habit formed in life may speak at death. ``The second is, to select in each month a day, on which you shall be most free from other employments, and to consecrate the entire of it to the great affair of learning to die well.''\endnote{Leonard of Port Maurice, ``On the Great Value Which a Religious Should Set on Making a Day's Retreat Each Month, for the Purpose of Securing a Good Death,'' from \emph{The Sacred Manual}, appended to Alphonsus Liguori, \emph{Preparation for Death} (Dublin: Duffy, 1843).} He wrote for religious, and his directions assume a cell; but the substance of the day can be kept by anyone who will set aside its hours.

The day is ordered toward the deathbed. ``The confession ought to be made as if it were the last of your life,'' with sorrow ``imagining that you must soon appear before the tribunal of God.'' The Communion is to be received ``as if you were going to receive the viaticum.''\endnote{Leonard, ``Day's Retreat,'' in \emph{Preparation for Death} (Duffy, 1843).} The morning meditation sets the soul in the room where it will die: ``you shall represent to yourself in the most vivid manner, the state in which you shall find yourself at the hour of death,---given over by the physicians, separated for ever from your companions, and admonished by your confessor to breathe your soul into the hands of God.'' It asks what one would then wish to have done toward God, toward oneself, and toward one's neighbour, and it ends in contrition and a purpose to live ``with as great fervour as if you expected death every day.''\endnote{Leonard, ``Day's Retreat,'' Duffy 1843.}

The evening meditation turns from fear to acceptance. Leonard gives five motives for accepting death willingly: its necessity, which the Son of God himself did not refuse; justice, since a sinner has deserved it; humility; love, because death opens the country ``where my sole occupation shall be to love you''; and resignation to the will of God, ``for its execution in the year, in the month, on the day, in the moment, in the form and manner which may be most pleasing to him.''\endnote{Leonard, ``Day's Retreat,'' Duffy 1843, with Psalm 4:9.} The day closes with an act of acceptance of death, and Leonard promises two fruits, the first on the word of the private revelations of Saint Gertrude. If death comes too suddenly for these acts to be made at the end, ``the Lord shall, according to St. Gertrude, accept them as if you actually performed them at the moment of your going to eternity''; and if there is time, the soul ``will have acquired by habit a facility of performing these acts.'' ``The man who has learned well in private the piece which he is to recite in public, delivers it on the stage with great ease.''\endnote{Leonard, ``Day's Retreat,'' Duffy 1843.}

\subsection{The protestation of death}

Alphonsus wrote such an act for his readers. He calls it a protestation of death, and it begins: ``My God, prostrate in your presence, I adore you; and I intend to make the following protestation, as if I were on the point of passing from this life to eternity.''\endnote{Alphonsus Liguori, ``Protestation of Death,'' in \emph{Preparation for Death}, Duffy 1843.} It professes the faith in which the Christian wishes to die. It hopes for pardon ``not through my merits, but through the merits of Jesus Christ,'' and it forestalls the last temptation: ``should the devil, at death, tempt me to despair at the sight of my sins, I protest that I wish always to hope in you, my Lord, and that I wish to die in the loving arms of your goodness.'' It forgives every enemy. It accepts death and its pains, ``and all the pains which shall accompany my death,'' in union with the death of Christ. It asks the Blessed Virgin, Saint Joseph, Saint Michael, and the guardian angel to stand by in ``that last battle.'' And it gives heart and soul to Jesus and Mary in advance, for the moment of the last breath: ``Jesus and Mary, I give you my heart and my soul.''\endnote{``Protestation of Death,'' Duffy 1843.}

With the protestation the same pages give a short prayer to be said every day for a good death:

\begin{quote}
O Lord Jesus Christ, through that agony which your most noble soul suffered when it went forth from your blessed body, have mercy on my sinful soul when it shall depart from my body. Amen.\endnote{``Prayer to be said every day for a good death,'' in \emph{Preparation for Death}, Duffy 1843.}
\end{quote}

\subsection{Watch ye, therefore}

The Church opens the new liturgical year with a Gospel of the end, as she closed the old one. On the first Sunday of Advent, in the 1962 Missal, she reads the signs in the sun and the moon and the stars, and the Lord's promise that ``this generation shall not pass away till all things be fulfilled.''\endnote{Luke 21:25--33, the Gospel of the First Sunday of Advent in the 1962 Missal; the sentence quoted is 21:32, Douay--Rheims (Challoner).} In the verses that follow, Saint Luke gives Christ's counsel: ``And take heed to yourselves, lest perhaps your hearts be overcharged with surfeiting and drunkenness and the cares of this life: and that day come upon you suddenly. \ldots{} Watch ye, therefore, praying at all times, that you may be accounted worthy to escape all these things that are to come and to stand before the Son of man.''\endnote{Luke 21:34, 36, Douay--Rheims (Challoner).}

Each of these practices is a way of keeping that command to watch. A Christian who keeps a day each month for his death, who chooses by the light of his last hour, who confesses as though for the last time and receives Communion as though it were his Viaticum, who says each day the prayer for a good death and each Hail Mary with attention to its last words, has done what Ecclesiasticus asks: in all his works he has remembered his last end. He will have met the fear of death as the martyrs met it, by loving Christ more than he fears to die. When the hour comes, whether with warning or without, the Church will be there with her sacraments and her prayers, and he will know her words, because he has said them all his life.
